\label{sec:disc}

\subsection{JKP Factor Returns: A Structural Mismatch, Not a Failure}

JKP factor returns are market-level, time-series signals: they take the same value for all stocks within a given month, encoding information about recent factor-portfolio performance. This structure is well suited to time-series allocation or regime identification, but not to firm-level cross-sectional stock selection.

In a long-short portfolio, returns are measured as long minus short \textit{within} the same month. A prediction that is constant across stocks in a given month produces a random within-month ranking, yielding zero expected return before costs and negative net returns after. This explains why JKP-only yields undefined monthly Rank IC and $-4.5\%$ net annualised return: not because JKP factors lack information, but because that information is aligned with time-series variation rather than firm-level cross-sectional differences.

The \textit{global} Rank IC (Spearman across all stock-month pairs) can appear positive for JKP models because it captures cross-month variation: correctly calling ``good'' versus ``bad'' months inflates global IC even when within-month stock rankings are uninformative. The combination paradox (JKP + FinBERT has higher global IC than FinBERT alone yet worse portfolio performance) follows directly: added JKP features capture time-series variation that inflates global IC but degrades cross-sectional discrimination, while also raising portfolio turnover from 58\% to 91\% ($\approx$3\% additional annual transaction costs). The same pattern appears across both XGBoost and deep learning models. A more natural cross-sectional benchmark would use firm-level characteristics (e.g.\ momentum, size, book-to-market) that vary across stocks within each month; we did not test such a benchmark, and future work using Compustat or CRSP-based predictors would provide a tighter comparison.

\subsection{XGBoost vs.\ Deep Learning}

\textbf{XGBoost vs.\ deep learning.} On the same 12 FinBERT features, XGBoost achieves monthly IC 0.032 and Sharpe 0.79, versus IC 0.025 and Sharpe 0.58 for MLP\textsubscript{FinBERT}. Both generate positive cross-sectional IC and profitable net returns, so the deep learning component delivers genuine empirical value rather than satisfying a formal requirement. XGBoost performs best, likely because tree-based splits handle the irregular sparsity of sentiment features efficiently; this is an empirical finding specific to this setup. The Ridge linear benchmark (IC 0.002, Sharpe $-0.10$) performs poorly, suggesting the signal requires non-linear modelling.

\textbf{Sequence models do not generalise across regimes.} LSTM\textsubscript{FinBERT} and Attention\textsubscript{FinBERT} model the temporal trajectory of a firm's earnings calls over up to 8 quarters. Both models perform well in 2019--2020 but deteriorate sharply in 2021--2023 (see Appendix~E). The quarterly trajectory patterns learned on 2008--2018 data do not transfer to the post-COVID rate-hike regime. Sequence models are less robust because quarterly earnings-call trajectories are sparse and regime-dependent; the cross-sectional rank of current sentiment is more regime-stable.

\textbf{Economic interpretation.} The results suggest that what matters is \textit{how a firm's tone ranks against peers today}, not \textit{how its tone has evolved over the past two years}. This is consistent with findings in the related analyst-report literature \cite{huang2023}, where recent tone changes carry more predictive content than multi-period sentiment histories.

\subsection{Limitations}

\begin{enumerate}[noitemsep]
    \item \textbf{Universe:} FinBERT models cover only earnings-call firms ($\approx$mid- and large-cap). Results may not generalise to small-cap or infrequently covered stocks.
    \item \textbf{Statistical power:} five years (60 months) of untouched test data leave meaningful uncertainty around the Sharpe estimate; the 108-month post-training IC path ($t = 3.79$) is a useful persistence check, but 2015--2018 is also used for early stopping.
    \item \textbf{Freshness cutoff:} the 3-month threshold was set a priori as a compromise between signal freshness and cross-sectional coverage. Appendix~C shows the 1-month cutoff yields $3\times$ higher global IC, but it sharply reduces the investable universe and may produce less stable decile portfolios. Optimising the cutoff ex post on backtest performance would risk data mining.
    \item \textbf{Transaction costs:} 0.2\% one-way is conservative for large-caps but may underestimate costs for smaller names. Sensitivity analysis (Appendix~D) shows FinBERT-only remains profitable up to 0.4\% one-way.
    \item \textbf{JKP comparison:} JKP-only is evaluated on a broader universe (40,872 test obs.) since it does not require earnings calls; FinBERT models require a recent call ($\leq$3 months) and, for delta features, at least two consecutive calls (36,901 test obs.). This partially confounds cross-model comparisons.
    \item \textbf{FinBERT context window:} the 512-token limit requires chunking; very long Q\&A sections may lose information in later chunks.
    \item \textbf{Sector neutrality:} long-short decile portfolios are not sector-neutral. Part of the signal may reflect sector-level earnings-season timing or industry-wide sentiment rather than firm-specific information. Checking sector concentration in the long and short deciles, or applying GICS/NAICS neutralisation, is a natural extension.
    \item \textbf{Implementation:} the backtest does not impose liquidity filters, market-cap weighting, borrow fees, or short-sale constraints. The mid- and large-cap tilt of the earnings-call universe improves practical implementability, but these extensions are necessary before any live deployment.
\end{enumerate}
